% Source: https://artofproblemsolving.com/community/c3h2956548_hexagon__chord
% Stable identity: aops-topic-2956548-q01
Let each side of length $3$ subtend a central angle $2\alpha$ and each side of length $5$ a central angle $2\beta$. Then $6\alpha+6\beta=360^\circ$, so $\beta=60^\circ-\alpha$. With radius $R$,
\[3=2R\sin\alpha,\qquad 5=2R\sin(60^\circ-\alpha).\]
\begin{asy}
usepackage("amsmath");
size(4.5cm);
defaultpen(fontsize(10pt));
real R=7/sqrt(3);real a=2*aTan(3*sqrt(3)/13),bb=120-a;
pair[] P;real t=90+1.5*a;for(int k=0;k<6;++k){P.push(R*dir(t));t-=(k<3?a:bb);}
draw(circle((0,0),R),gray);draw(P[0]--P[1]--P[2]--P[3]--P[4]--P[5]--cycle);draw(P[0]--P[3],linewidth(1));
label("$3$",(P[0]+P[1])/2,NE);label("$3$",(P[1]+P[2])/2,E);label("$3$",(P[2]+P[3])/2,NE);label("$5$",(P[3]+P[4])/2,S);label("$5$",(P[4]+P[5])/2,W);label("$5$",(P[5]+P[0])/2,NW);dot((0,0));
\end{asy}

Therefore $3\left(\dfrac{\sqrt3}2\cos\alpha-\dfrac12\sin\alpha\right)=5\sin\alpha$, so $\tan\alpha=\dfrac{3\sqrt3}{13}$. This gives $\sin\alpha=\dfrac{3\sqrt3}{14}$ and $R=\dfrac{3}{2\sin\alpha}=\dfrac7{\sqrt3}$.

The chord spans the three short sides, so its length is
\[2R\sin3\alpha=2R(3\sin\alpha-4\sin^3\alpha)=\dfrac{14}{\sqrt3}\cdot\dfrac{180\sqrt3}{343}=\dfrac{360}{49}.\]
So $m+n=\boxed{409}$.
